% Vietnamese translation for OI Public Library
% Source-PDF: 其他 OTHERS/猜数字游戏的计算机求解 - 莫涛.pdf
\documentclass[11pt]{article}
\usepackage{oipl-vn}

\newcommand{\AB}[2]{#1A#2B}
\newcommand{\SetS}{\mathcal{S}}
\newcommand{\Ans}{\mathcal{A}}

\title{Giải trò chơi đoán số bằng máy tính}
\author{Mo Tao\\Khoa Máy tính, Đại học Thanh Hoa}
\date{Tháng 1 năm 2011}

\begin{document}
\maketitle

\origtitle{Giải trò chơi đoán số bằng máy tính}
\sourcepath{Others / Bulls and Cows computer solving - Mo Tao}

\begin{abstract}
Bài viết khảo sát nhiều chiến lược để máy tính giải trò chơi đoán số
\term{Bulls and Cows}. Trọng tâm là cách giảm số lần đoán trung bình, tăng tốc
các thuật toán heuristic, và hiện thực những thuật toán đó bằng C++. Bản dịch
giữ cấu trúc bài gốc: luật chuẩn và biến thể, phương pháp lọc tập ứng viên,
các chiến lược đơn giản, các hàm đánh giá heuristic, tối ưu bằng lớp tương
đương, bảng kết quả thực nghiệm, cùng mô tả thư viện kiểm thử và chương trình
chính.
\end{abstract}

\paragraph{Từ khóa.}
Phương pháp lọc; hàm đánh giá; tối ưu bằng lớp tương đương; Bulls and Cows;
Mastermind.

\tableofcontents

\section{Mở đầu}

Trò chơi đoán số, còn gọi là \term{Bulls and Cows}, là một trò chơi trí tuệ
nhỏ được cho là phổ biến ở Anh từ khoảng giữa thế kỷ 20. Trò chơi thường do
hai người chơi, cũng có thể là một người chơi với máy tính, và có thể chơi
trên giấy hoặc trực tuyến. Luật chơi đơn giản, nhưng lại kiểm tra tính chặt
chẽ và sự kiên nhẫn của người chơi.

\subsection{Luật chuẩn}

Thông thường một người ra số và một người đoán. Người ra số nghĩ trước một số
gồm bốn chữ số không lặp lại, trong đó chữ số đầu tiên được phép là \(0\), và
không cho người đoán biết. Mỗi lần người đoán đưa ra một số, người ra số trả
lời theo dạng \(xAyB\): \(x\) là số chữ số đúng và đúng vị trí, còn \(y\) là
số chữ số đúng nhưng sai vị trí.

Ví dụ, nếu đáp án là \(5234\) và người đoán đưa ra \(5346\), phản hồi là
\(\AB{1}{2}\). Chữ số \(5\) đúng vị trí nên tính \(1A\); hai chữ số \(3\) và
\(4\) có trong đáp án nhưng sai vị trí nên tính \(2B\). Người đoán tiếp tục
dựa vào các phản hồi đó cho đến khi đoán trúng, tức nhận \(\AB{4}{0}\).

Trò chơi thường có giới hạn số lần đoán. Theo tính toán bằng máy tính, nếu
dùng chiến lược chặt chẽ thì với luật chuẩn mọi đáp án đều có thể đoán ra
trong nhiều nhất \(7\) lần. Cần chú ý rằng một số nơi định nghĩa giới hạn là
``sau bao nhiêu lần đoán thì đã có thể xác định chắc chắn đáp án''. Theo cách
đó, đôi khi người chơi vẫn cần thêm một lần đoán nữa để thật sự nhận phản hồi
\(\AB{4}{0}\).

\subsection{Các biến thể}

Luật chuẩn dùng \(10\) chữ số \(0\) đến \(9\) và \(4\) vị trí. Có thể thay đổi
số chữ số hoặc số vị trí để tạo biến thể, chẳng hạn dùng \(9\) chữ số cho trò
chơi \(4\) vị trí.

Một biến thể khác cho phép lặp chữ số; trò chơi này thường được gọi là
\term{Mastermind}. Ngoài luật chuẩn, khi có chữ số lặp, mỗi lần xuất hiện chỉ
được ghép một lần và lấy kết quả tốt nhất. Ví dụ đáp án là \(5543\), đoán
\(5255\). Không thể ghép chữ số \(5\) đầu tiên của lượt đoán với chữ số \(5\)
thứ hai của đáp án nếu điều đó làm kết quả kém tối ưu. Kết quả đúng là
\(\AB{1}{1}\): chữ số \(5\) đầu đúng vị trí, và một trong hai chữ số \(5\) còn
lại của lượt đoán ghép với chữ số \(5\) thứ hai của đáp án. Nếu đoán \(5267\)
thì chỉ có \(\AB{1}{0}\), vì chữ số \(5\) đầu đã được tính đúng vị trí và
không thể dùng lại để ghép sai vị trí.

Mastermind chuẩn có \(4\) vị trí, mỗi vị trí lấy một trong \(6\) giá trị, khi
chơi thường viết bằng \(A\) đến \(F\). Các biến thể khác gồm Royale
Mastermind với \(5\) vị trí và mỗi vị trí lấy một trong \(5\) giá trị,
Mastermind Challenge trong đó hai bên thay phiên ra đề cho nhau, và Word
Mastermind trong đó \(0\) đến \(25\) được xem là \(A\) đến \(Z\), còn chuỗi
cần đoán phải là một từ.

Một biến thể khó hơn là \term{Static Game}: các lượt đoán không được phụ thuộc
vào phản hồi. Người chơi đưa ra một lần một tập các lượt đoán; với mọi đáp án,
các phản hồi tương ứng phải đủ để xác định đáp án. Với Mastermind chuẩn,
Greenwell đã đưa ra một nghiệm chỉ gồm \(6\) lượt đoán:
\[
\begin{gathered}
(1,2,2,1),\ (2,3,5,4),\ (3,3,1,1),\\
(4,5,2,4),\ (5,6,5,6),\ (6,6,4,3).
\end{gathered}
\]
Về mặt lý thuyết \(5\) lượt đoán là đủ, nhưng bài gốc ghi nhận rằng khi đó
chưa tìm được nghiệm như vậy.

\section{Chiến lược giải với luật chuẩn}

Trò chơi đoán số là một bài toán tương tác thông tin cổ điển: người chơi cần
dùng càng ít câu hỏi càng tốt để thu thập thông tin từ người ra số, rồi suy ra
đáp án. Điểm phức tạp là các câu hỏi được đặt tuần tự; người chơi có thể dựa
vào phản hồi trước đó để điều chỉnh chiến lược.

Phần này trước hết thống nhất tiêu chuẩn đánh giá, sau đó trình bày phương
pháp lọc, là thuật toán phổ biến nhất cho trò chơi này. Các chiến lược đơn
giản và heuristic sẽ được so sánh về cơ sở lý thuyết, hiện thực, kết quả chạy
và độ phức tạp. Cuối cùng là một kỹ thuật tối ưu quan trọng cho thuật toán
heuristic và phần mô tả mã nguồn.

\subsection{Tiêu chuẩn thống nhất}

Để so sánh các chiến lược, bài gốc dùng các quy ước sau:
\begin{enumerate}
  \item Chơi theo luật chuẩn: \(4\) vị trí, mỗi vị trí lấy một chữ số trong
  \(0\) đến \(9\), và các chữ số khác nhau.
  \item Trò chơi kết thúc khi một lượt đoán nhận \(\AB{4}{0}\); lượt đoán này
  cũng được tính vào tổng số lần đoán.
  \item Người ra số công bằng: không rò rỉ thông tin, không trả lời sai, và
  không đổi đáp án giữa chừng để gây khó cho người đoán.
  \item Mọi ván phải kết thúc trong nhiều nhất \(7\) lần đoán; dưới ràng buộc
  đó, giảm số lần đoán trung bình. Kết quả được thống kê bằng cách thử toàn bộ
  \(5040\) đáp án.
  \item Tối ưu thời gian chạy; mỗi ván nên chạy không quá \(1\) giây.
\end{enumerate}

Với quy ước thứ hai, một chiến lược có thể tình cờ đoán trúng khi chưa chắc
chắn, hoặc ngược lại đã xác định được đáp án nhưng vẫn phải dùng thêm một lượt
đoán để nhận \(\AB{4}{0}\). Về kỳ vọng, chi tiết này không làm thay đổi bản
chất, và nó phù hợp hơn với trực giác rằng trò chơi kết thúc khi ``đoán
trúng'', không chỉ khi ``biết đáp án''.

Quy ước thứ ba loại bỏ người ra số đối kháng. Nếu người ra số được phép chọn
đáp án thích nghi nhưng vẫn giữ cho phản hồi nhất quán, họ có thể làm chiến
lược của người đoán tệ đi. Ví dụ, khi người đoán đã xác định ba vị trí đầu là
\(012\) và định lần lượt thử \(0123,0124,\ldots,0129\), trung bình chỉ cần
\(7/2=3.5\) lượt; nhưng một người ra số đối kháng có thể chọn \(0129\), buộc
người đoán dùng \(7\) lượt. Vì vậy bài gốc giả định đáp án được cố định ngay
từ đầu, hoặc được quản lý bởi một chương trình thư viện.

Quy ước thứ tư phản ánh hai mục tiêu thường gặp: giảm số lần đoán tệ nhất và
giảm số lần đoán trung bình. Trong một số biến thể, hai mục tiêu không đồng
thời đạt được. Chẳng hạn với Mastermind, chiến lược có trung bình tốt nhất cần
khoảng \(4.340\) lượt nhưng trường hợp xấu nhất cần \(6\) lượt; nếu buộc nhiều
nhất \(5\) lượt, trung bình tốt nhất tăng lên khoảng \(4.341\). Với luật chuẩn
Bulls and Cows, hầu hết chiến lược trung bình tốt cũng chỉ cần \(7\) lượt
trong trường hợp xấu nhất, nên ràng buộc này là hợp lý.

Với quy ước thứ năm, nếu bỏ giới hạn thời gian thì máy tính có thể duyệt toàn
bộ cây trò chơi hữu hạn để tìm chiến lược tối ưu. Tuy nhiên, ngay cả khi mỗi
ván chỉ mất \(1\) giây, thử hết \(5040\) đáp án cũng mất \(5040\) giây, tức
hơn \(1\) giờ \(24\) phút. Do đó tối ưu thời gian chạy cũng trực tiếp giúp
thử nghiệm và điều chỉnh chiến lược tốt hơn. Bài gốc đưa ra một tối ưu giúp
giảm thời gian chạy từ hơn \(3000\) giây xuống dưới \(100\) giây.

\subsection{Phương pháp lọc}

Với luật chuẩn, số đáp án chỉ là
\[
  10\cdot 9\cdot 8\cdot 7 = 5040.
\]
Với Mastermind chuẩn là \(6^4=1296\), với Royale Mastermind là \(5^5=3125\).
Số khả năng nhỏ là một đặc điểm tự nhiên của trò chơi: nếu không, người chơi
người thật khó có thể suy luận thủ công.

Một đặc điểm khác là miền giá trị ở mỗi vị trí không quá lớn. Nếu mỗi vị trí
có thể là \(0\) đến \(99\), chỉ \(2\) vị trí đã có \(100^2=10000\) khả năng.
Trực giác khi con người giải trò chơi này thường là trước hết xác định các chữ
số xuất hiện trong đáp án, cố gắng nhận phản hồi kiểu \(\AB{1}{3}\) hoặc
\(\AB{2}{2}\), rồi hoán vị chúng để tìm đúng vị trí. Miền chữ số quá rộng làm
cách suy luận này kém hiệu quả.

Có hai hướng thiết kế thuật toán. Hướng thứ nhất mô phỏng kỹ thuật của con
người bằng một quy trình đoán và suy luận logic. Cách này khó lập trình và kết
quả không ổn định, vì nó chưa khai thác trực tiếp đặc điểm bản chất: tập kết
quả có thể là hữu hạn và nhỏ. Hướng thứ hai là \term{phương pháp lọc}, được
dùng xuyên suốt bài này.

Thuật toán lọc:
\begin{enumerate}
  \item Khởi tạo tập ứng viên \(\SetS\) gồm toàn bộ \(5040\) đáp án có thể.
  \item Dùng một chiến lược nào đó để chọn lượt đoán \(X\), nhận phản hồi
  \(R=xAyB\).
  \item Xóa khỏi \(\SetS\) mọi đáp án \(Y\) mà nếu đáp án là \(Y\) thì đoán
  \(X\) không cho phản hồi \(R\).
  \item Nếu \(\SetS\) còn hơn một đáp án, quay lại bước 2.
  \item Ứng viên duy nhất còn lại là đáp án.
\end{enumerate}

Mã giả:
\begin{verbatim}
Solve()
    S <- all possible answers
    while |S| > 1 do
        X <- find(S)
        R <- feedback for guess X
        for every Y in S do
            if feedback(answer = Y, guess = X) != R then
                remove Y from S
            end if
        end for
    end while
    return the only element of S
end Solve
\end{verbatim}

Tên ``lọc'' đến từ việc mỗi phản hồi loại bỏ các đáp án không còn nhất quán.
Bước then chốt là hàm \texttt{find}: từ tập ứng viên hiện tại, chọn lượt đoán
sao cho trung bình cần ít lượt đoán nhất.

\subsection{Các chiến lược đơn giản}

Một ý tưởng trực quan là chỉ đoán các phần tử trong tập ứng viên. Như vậy bất
kể phản hồi ra sao, lượt đoán vẫn có xác suất trúng đáp án. Câu hỏi còn lại
là chọn phần tử nào trong \(\SetS\).

Chiến lược thứ nhất sắp \(\SetS\) theo thứ tự từ điển và chọn phần tử nhỏ
nhất:
\begin{verbatim}
find(S)
    sort S in lexicographic order
    return first element of S
end find
\end{verbatim}
Với hai chuỗi số khác nhau \(A,B\), tìm vị trí nhỏ nhất \(p\) sao cho
\(A_p\ne B_p\); nếu \(A_p<B_p\) thì \(A\) nhỏ hơn \(B\) theo thứ tự từ điển.

Kết quả:
\begin{center}
\small
\begin{tabular}{|l|r|}
\hline
Số lần đoán trung bình & 5.560 \\
Số lần đoán nhiều nhất & 9 \\
Tổng thời gian & 10.367 giây \\
\hline
\end{tabular}

\medskip
\begin{tabular}{|l|rrrrrrrrr|}
\hline
Số lần & 1 & 2 & 3 & 4 & 5 & 6 & 7 & 8 & 9 \\
\hline
Số ván & 1 & 13 & 108 & 596 & 1668 & 1768 & 752 & 129 & 5 \\
\hline
\end{tabular}
\end{center}

Chiến lược thứ hai chọn ngẫu nhiên đều một phần tử của \(\SetS\):
\begin{verbatim}
find(S)
    return a uniformly random element of S
end find
\end{verbatim}

Kết quả:
\begin{center}
\small
\begin{tabular}{|l|r|}
\hline
Số lần đoán trung bình & 5.465 \\
Số lần đoán nhiều nhất & 9 \\
Tổng thời gian & 10.508 giây \\
\hline
\end{tabular}

\medskip
\begin{tabular}{|l|rrrrrrrrr|}
\hline
Số lần & 1 & 2 & 3 & 4 & 5 & 6 & 7 & 8 & 9 \\
\hline
Số ván & 1 & 10 & 114 & 593 & 1844 & 1804 & 626 & 47 & 1 \\
\hline
\end{tabular}
\end{center}

Chiến lược ngẫu nhiên tốt hơn rõ rệt: số ván cần hơn \(7\) lượt giảm mạnh.
Dù vậy, cả hai đều không thỏa quy ước ``nhiều nhất \(7\) lượt''. Ưu điểm của
chúng là thời gian chạy chỉ khoảng \(10\) giây cho toàn bộ thử nghiệm.

\subsection{Chiến lược heuristic}

Chiến lược đơn giản kém vì lượt đoán còn quá tùy ý. Ta có thể duyệt mọi chuỗi
số hợp lệ làm lượt đoán, rồi chọn chuỗi ``tốt nhất''.

Vấn đề thứ nhất là định nghĩa ``tốt nhất''. Duyệt toàn bộ phần còn lại của cây
trò chơi để chọn lượt có số lần đoán nhỏ nhất là lý tưởng, nhưng không khả thi
về thời gian. Vì vậy ta dùng \term{hàm đánh giá}: một hàm dùng thông tin hữu
hạn hiện có để ước lượng chất lượng của một hành động và biểu diễn chất lượng
đó bằng một con số. Trong trò chơi đoán số, một đánh giá tự nhiên là kích
thước tập ứng viên còn lại sau lượt đoán: tập càng nhỏ thì thường cần ít lượt
đoán hơn.

Vấn đề thứ hai là có nên chỉ duyệt các phần tử của \(\SetS\) hay không. Với
mục tiêu thu thập thông tin, một chuỗi không thể là đáp án vẫn có thể là lượt
đoán tốt hơn. Nếu hàm đánh giá đủ hợp lý, miền duyệt rộng hơn thường cho kết
quả tốt hơn. Bài gốc dùng cách dung hòa: duyệt mọi đáp án hợp lệ trong
\(\Ans\), và khi giá trị đánh giá bằng nhau thì ưu tiên phần tử thuộc
\(\SetS\).

\begin{verbatim}
find(S)
    ret <- 0123
    for every valid answer X in A do
        if value(X) < value(ret) then ret <- X
        if value(X) = value(ret) and X is in S then ret <- X
    end for
    return ret
end find
\end{verbatim}

Quy ước là \(\texttt{value}(X)\) càng nhỏ thì lượt đoán \(X\) càng tốt.

\subsection{Thiết kế hàm đánh giá}

\subsubsection{Kích thước tập ứng viên tệ nhất}

Hàm đánh giá chỉ dựa trên tập ứng viên hiện tại \(\SetS\) và lượt đoán đang
xét \(X\). Với một lượt đoán \(X\), phản hồi có \(15\) loại:
\[
\begin{gathered}
\AB{0}{0},\AB{0}{1},\AB{0}{2},\AB{0}{3},\AB{0}{4},\\
\AB{1}{0},\AB{1}{1},\AB{1}{2},\AB{1}{3},\\
\AB{2}{0},\AB{2}{1},\AB{2}{2},\\
\AB{3}{0},\AB{3}{1},\AB{4}{0}.
\end{gathered}
\]
Gọi \(\SetS[a][b]\) là tập ứng viên còn lại nếu đoán \(X\) và nhận phản hồi
\(aAbB\). Hàm đánh giá theo kích thước tệ nhất là:
\begin{verbatim}
value(X)
    compute S[a][b]
    ret <- 0
    for a = 0..4 do
        for b = 0..4-a do
            ret <- max(ret, |S[a][b]|)
        end for
    end for
    return ret
end value
\end{verbatim}

Kết quả:
\begin{center}
\small
\begin{tabular}{|l|r|}
\hline
Số lần đoán trung bình & 5.389 \\
Số lần đoán nhiều nhất & 7 \\
Tổng thời gian & 79.620 giây \\
\hline
\end{tabular}

\medskip
\begin{tabular}{|l|rrrrrrrrr|}
\hline
Số lần & 1 & 2 & 3 & 4 & 5 & 6 & 7 & 8 & 9 \\
\hline
Số ván & 1 & 3 & 44 & 519 & 2106 & 2152 & 215 & 0 & 0 \\
\hline
\end{tabular}
\end{center}

Chiến lược này thỏa ràng buộc \(7\) lượt. Trung bình giảm từ \(5.465\) xuống
\(5.389\), tức trong \(5040\) ván, xấp xỉ cứ \(10\) ván thì tiết kiệm được một
lượt so với chiến lược đơn giản. Đây là cải thiện đáng kể, đồng thời cho thấy
trò chơi khá khó: sau khi đã có thuật toán hợp lý, cải thiện nhỏ cũng không
dễ.

Nếu hiện thực ngây thơ, với mỗi \(X\) ta quét mọi phần tử của \(\SetS\), tính
phản hồi và phân loại vào \(\SetS[a][b]\). Một quyết định có độ phức tạp
\(O(5040\cdot |\SetS|\cdot |R|)\). Thử mọi ván có độ phức tạp xấp xỉ
\(O(|\SetS|^3)\) với hằng số lớn. Chạy trực tiếp mất hơn một giờ; các thời
gian trong bài đã dùng tối ưu ở mục sau.

\subsubsection{Kích thước tập ứng viên trung bình}

Thay vì tối thiểu hóa trường hợp xấu nhất, ta có thể tối thiểu hóa kích thước
kỳ vọng của tập còn lại. Xét ví dụ chỉ có hai vị trí, tại một thời điểm
\(\SetS=\{12,21,34\}\). Nếu đoán \(34\), đáp án \(34\) cho phản hồi
\(\AB{2}{0}\), còn \(12\) và \(21\) cho \(\AB{0}{0}\). Kích thước còn lại kỳ
vọng là
\[
  1\cdot \frac13 + 2\cdot \frac23 = \frac53.
\]
Nếu đoán \(12\), ba đáp án lần lượt cho ba phản hồi khác nhau
\(\AB{2}{0},\AB{0}{2},\AB{0}{0}\), nên kích thước còn lại kỳ vọng là \(1\).
Vì vậy đoán \(12\) tốt hơn.

Hàm đánh giá:
\begin{verbatim}
value(X)
    compute S[a][b]
    ret <- 0
    for a = 0..4 do
        for b = 0..4-a do
            ret <- ret + |S[a][b]|^2
        end for
    end for
    return ret
end value
\end{verbatim}
Giá trị kỳ vọng thật sự còn chia cho \(|\SetS|\), nhưng trong cùng một quyết
định mẫu số không đổi nên có thể bỏ qua.

Kết quả:
\begin{center}
\small
\begin{tabular}{|l|r|}
\hline
Số lần đoán trung bình & 5.268 \\
Số lần đoán nhiều nhất & 7 \\
Tổng thời gian & 73.762 giây \\
\hline
\end{tabular}

\medskip
\begin{tabular}{|l|rrrrrrrrr|}
\hline
Số lần & 1 & 2 & 3 & 4 & 5 & 6 & 7 & 8 & 9 \\
\hline
Số ván & 1 & 4 & 59 & 573 & 2430 & 1886 & 87 & 0 & 0 \\
\hline
\end{tabular}
\end{center}

Trung bình tiếp tục cải thiện rõ rệt, trong khi các số liệu khác gần như không
đổi. Hàm đánh giá này khớp tốt hơn với chất lượng thực tế của lượt đoán.

\subsubsection{Số bước kỳ vọng}

Một hướng tự nhiên hơn là dùng kỳ vọng số lượt đoán còn lại:
\begin{verbatim}
value(X)
    compute S[a][b]
    ret <- 0
    for a = 0..4 do
        for b = 0..4-a do
            ret <- ret + |S[a][b]| * Step(S[a][b])
        end for
    end for
    return ret
end value
\end{verbatim}
Ở đây \(\operatorname{Step}(\SetS[a][b])\) là số bước kỳ vọng để giải tiếp từ
tập đó. Vì không thể tính chính xác trong giới hạn thời gian, bài gốc ước
lượng rằng mỗi lượt làm tập ứng viên thu nhỏ theo một tỷ lệ trung bình \(k\),
tức
\[
  \operatorname{Step}(\SetS[a][b]) = \log_k |\SetS[a][b]|.
\]
Giá trị cụ thể của \(k\) không ảnh hưởng đến thứ tự so sánh, nên có thể dùng
\[
  \operatorname{Step}(\SetS[a][b]) = \ln |\SetS[a][b]|.
\]

Kết quả:
\begin{center}
\small
\begin{tabular}{|l|r|}
\hline
Số lần đoán trung bình & 5.265 \\
Số lần đoán nhiều nhất & 7 \\
Tổng thời gian & 95.682 giây \\
\hline
\end{tabular}

\medskip
\begin{tabular}{|l|rrrrrrrrr|}
\hline
Số lần & 1 & 2 & 3 & 4 & 5 & 6 & 7 & 8 & 9 \\
\hline
Số ván & 1 & 4 & 53 & 560 & 2515 & 1796 & 111 & 0 & 0 \\
\hline
\end{tabular}
\end{center}

So với hàm kích thước trung bình, kết quả không khác nhiều. Tuy nhiên cách
nhìn qua hàm \(\operatorname{Step}\) có khả năng mở rộng tốt hơn và sẽ xuất
hiện lại trong phần chiến lược tối ưu hóa.

\subsubsection{Số loại phản hồi}

Kooi đề xuất năm 2005 một chiến lược dựa trên ý tưởng sau: tập còn lại nhỏ hơn
thì thuận lợi hơn cho những lượt sau, nhưng xác suất nó xảy ra cũng nhỏ hơn;
vì vậy có thể xem đóng góp của mỗi nhánh phản hồi là như nhau. Do đó với một
lượt đoán \(X\), càng có nhiều tập còn lại khác rỗng, tức càng có nhiều loại
phản hồi có thể xuất hiện, thì \(X\) càng tốt.

\begin{verbatim}
value(X)
    compute S[a][b]
    ret <- 0
    for a = 0..4 do
        for b = 0..4-a do
            if S[a][b] is not empty then ret <- ret + 1
        end for
    end for
    return -ret
end value
\end{verbatim}
Vì \(\texttt{value}\) được tối thiểu hóa còn số loại phản hồi cần tối đa hóa,
ta trả về giá trị đối.

Kết quả:
\begin{center}
\small
\begin{tabular}{|l|r|}
\hline
Số lần đoán trung bình & 5.308 \\
Số lần đoán nhiều nhất & 8 \\
Tổng thời gian & 78.632 giây \\
\hline
\end{tabular}

\medskip
\begin{tabular}{|l|rrrrrrrrr|}
\hline
Số lần & 1 & 2 & 3 & 4 & 5 & 6 & 7 & 8 & 9 \\
\hline
Số ván & 1 & 11 & 80 & 556 & 2277 & 1929 & 183 & 3 & 0 \\
\hline
\end{tabular}
\end{center}

Với luật chuẩn, hàm này kém hơn: có \(3\) ván cần \(8\) lượt và trung bình cao
hơn. Nhưng với Mastermind, đây là một trong các chiến lược tốt đã biết. Điều
này cho thấy cùng một thuật toán có thể khác biệt lớn khi đổi luật chơi.

\subsection{Tối ưu thuật toán}

Nút thắt là tính các tập \(\SetS[a][b]\). Thuật toán ngây thơ duyệt từng lượt
đoán \(X\), quét \(\SetS\), tính phản hồi và phân loại. Độ phức tạp khó giảm
về bậc, nhưng có thể giảm rất mạnh hằng số.

Trong một ván, thời gian chủ yếu nằm ở vài lượt đầu, khi \(\SetS\) còn lớn.
Về cuối ván, \(|\SetS|\) gần như là hằng số nhỏ. Ví dụ với hàm số bước kỳ
vọng, khi đáp án là \(9876\), một quá trình có thể là:
\begin{center}
\small
\begin{tabular}{|l|l|r|r|}
\hline
Lượt & Đoán và phản hồi & \(|\SetS|\) & Số \(X\) khác bản chất \\
\hline
1 & \(0123,\ \AB{0}{0}\) & 5040 & 1 \\
2 & \(4567,\ \AB{0}{2}\) & 360 & 209 \\
3 & \(8975,\ \AB{1}{2}\) & 84 & 3360 \\
4 & \(7948,\ \AB{0}{3}\) & 21 & 5040 \\
5 & \(9876,\ \AB{4}{0}\) & 5 & 5040 \\
\hline
\end{tabular}
\end{center}

Ở đầu ván, nhiều lượt đoán là tương đương về bản chất. Lượt đầu, mọi \(X\) hợp
lệ đều không khác nhau; có thể luôn đoán \(0123\) mà không ảnh hưởng đến số
lượt trung bình. Lượt thứ hai, do các chữ số \(4,5,6,7,8,9\) đều chưa từng
xuất hiện, các lượt như \(0142\) và \(0172\) cũng tương đương.

Bài gốc dùng quy tắc sinh đại diện lớp tương đương như sau. Khi duyệt \(X\),
các chữ số chưa từng dùng phải xuất hiện theo thứ tự tăng dần. Cụ thể, nếu
chữ số \(P\) chưa xuất hiện trong các lượt đoán trước, thì \(X\) chỉ được chứa
\(P\) khi \(P-1\) đã xuất hiện trong lượt đoán trước đó hoặc đã xuất hiện ở
phần đầu của chính \(X\). Nhờ vậy mỗi lớp tương đương chỉ được duyệt một lần.

\begin{verbatim}
Dfs(i)
    if i = 4 then
        output one essentially different X
    else
        for P = 0..9 do
            if P satisfies the representative rule then
                Dfs(i + 1)
            end if
        end for
    end if
end Dfs
\end{verbatim}

Hiệu quả rất lớn. Trong ví dụ trên, lượt tính nặng nhất là lượt thứ ba:
\[
  |\SetS|\cdot \#X = 84\cdot 3360 = 282240,
\]
nhỏ hơn rất nhiều so với ước lượng ngây thơ \(5040\cdot 5040=25401600\).

Một tối ưu hằng số khác nằm ở cách tính phản hồi. Nếu duyệt \(X\) ở vòng ngoài,
mỗi phần tử của \(\SetS\) lại phải dựng lại mảng đánh dấu chữ số để đếm \(B\).
Nếu lưu trước các \(X\) đại diện, rồi đặt vòng lặp trên \(\SetS\) ở ngoài và
quét các \(X\) ở trong, có thể tái sử dụng thông tin đánh dấu và giảm đáng kể
chi phí. Các chi tiết hiện thực khác cũng ảnh hưởng nhiều, nhưng bài gốc không
triển khai hết vì quá vụn. Sau các tối ưu này, thời gian chạy giảm xuống
khoảng \(100\) giây.

Về bản chất, phương pháp lọc heuristic là một tìm kiếm độ sâu \(1\). Các kỹ
thuật cắt tỉa thường dùng trong tìm kiếm đều có thể áp dụng để tối ưu nó.

\subsection{Mô tả mã nguồn và chương trình}

Các chương trình trong bài gốc được viết bằng C++ và chạy trên Windows 7 với
Dev-C++ 4.9.9.2. Thời gian thực nghiệm được đo trên máy cá nhân có RAM
\(2.00\) GB, CPU Intel Core i3 \(2.53\) GHz.

\subsubsection{Thư viện kiểm thử}

Bài gốc viết một lớp \texttt{game} để quản lý toàn bộ \(5040\) ván và chặn
truy cập trực tiếp vào đáp án, nhằm giữ công bằng. Cách dùng:
\begin{enumerate}
  \item Đặt mã vào \texttt{game.h} và thêm \verb|#include "game.h"| trong
  chương trình chính.
  \item Đổi biến \texttt{AIname} để đổi tên các tệp xuất log.
  \item Gọi \texttt{guess(a)} để đoán, trong đó \texttt{a} là mảng hoặc
  \texttt{vector} gồm \(4\) chữ số. Hàm trả về \texttt{pair<int,int>} chứa
  phản hồi \((A,B)\). Khi đoán trúng \(\AB{4}{0}\), thư viện tự chuyển sang
  ván tiếp theo. Nếu quá giới hạn mà chưa đoán được, ván được tính thất bại.
  Sau khi chạy hết \(5040\) ván, thư viện xuất log chi tiết và thống kê tổng.
\end{enumerate}

Lõi thư viện có các hằng số và giao diện sau:
\begin{verbatim}
const int L = 4, S = 10, maxt = 9, Size = 5040;
const pair<int, int> error = make_pair(-1, -1);
const pair<int, int> over = make_pair(-2, -2);

string AIname = "Noname";

class game {
    int P, Ps, T[maxt + 1];
    int a[L], t;
public:
    game();
    pair<int, int> guess(int* b);
};

pair<int, int> guess(int* b);
pair<int, int> guess(vector<int> b);
\end{verbatim}

Trong \texttt{guess}, thư viện kiểm tra chữ số có nằm trong miền hợp lệ và có
bị lặp hay không, tính số \(A\) bằng so sánh cùng vị trí, tính số chữ số chung
rồi trừ \(A\) để nhận \(B\). Khi một ván kết thúc, thư viện cập nhật bảng
\texttt{T[i]}: số ván đoán đúng sau \(i\) lượt. Khi hết mọi đáp án, thư viện
in số lần đoán trung bình, số lần nhiều nhất, tổng thời gian, và phân bố số
ván theo số lượt đoán.

\subsubsection{Chương trình chính}

Chương trình chính trong bài gốc dùng hàm đánh giá
\[
  \sum_{a,b} |\SetS[a][b]|\ln(|\SetS[a][b]|+1),
\]
tức biến thể tối ưu của hàm số bước kỳ vọng. Mã dùng \texttt{vector} để lưu
các chuỗi số. Những phần chính là:
\begin{itemize}
  \item \texttt{dfs}: sinh toàn bộ \(5040\) đáp án hợp lệ.
  \item \texttt{dfs2}: sinh các lượt đoán khác bản chất theo quy tắc lớp
  tương đương.
  \item \texttt{calc}: tính phản hồi nếu một chuỗi là đáp án và một chuỗi là
  lượt đoán.
  \item \texttt{find}: chọn lượt đoán tốt nhất bằng cách tính bảng
  \texttt{ct[i][j][k]}, trong đó \texttt{ct[i][j][k]} là kích thước tập ứng
  viên khi lượt đoán \texttt{h[i]} trả về \(jAkB\).
  \item \texttt{main}: lặp qua toàn bộ \(5040\) đáp án, duy trì tập ứng viên
  hiện tại, gọi \texttt{guess}, rồi lọc lại tập ứng viên bằng \texttt{calc}.
\end{itemize}

Đoạn lựa chọn hàm đánh giá có dạng:
\begin{verbatim}
double best = 1e100;
vector<int> ret;

for each candidate guess h[i] do
    double v = 0;
    for j = 0..L do
        for k = 0..L-j do
            v += ct[i][j][k] * log(ct[i][j][k] + 1);

            // Other evaluators used in the paper:
            // v += ct[i][j][k] * ct[i][j][k];
            // if (ct[i][j][k]) v += ct[i][j][k] * log(ct[i][j][k]);
            // v = max(v, double(ct[i][j][k]));
            // if (ct[i][j][k]) --v;
        end for
    end for

    if v < best, or v == best and h[i] is still possible, choose h[i]
end for
\end{verbatim}

Trong vòng chơi chính:
\begin{verbatim}
for every real answer do
    g <- all possible answers
    t <- 0
    p <- L - 1
    repeat
        t <- t + 1
        a <- find()
        ret <- guess(a)
        if ret.first == 4 then break

        update p by the largest digit used in a
        h <- empty
        for every y in g do
            if calc(y, a) == ret then append y to h
        end for
        g <- h
    forever
end for
\end{verbatim}

Các hàm đánh giá khác có thể kiểm thử bằng cách thay phần tính \texttt{v};
chiến lược đơn giản chỉ cần thay hàm \texttt{find}.

\section{Những thử nghiệm khác}

Trong quá trình viết bài, khi hiệu quả trung bình của heuristic gần chạm nút
thắt, tác giả tiếp tục suy nghĩ về cách cải thiện bản chất chiến lược và có
một số kết quả sơ bộ. Ngoài ra, dựa trên bài \texttt{guess} của Winter Camp
2006, tác giả cũng thử một số luật chơi không chuẩn.

\subsection{Chiến lược tối ưu hóa}

Chiến lược tối ưu đầy đủ khó tính trong thực tế, nhưng tối ưu hàm đánh giá
``số bước kỳ vọng'' có thể đưa ta đến khá gần. Trong mục trước,
\(\operatorname{Step}(\SetS[a][b])=\ln|\SetS[a][b]|\) khá hợp lý khi tập còn
lại lớn. Nhưng khi tập nhỏ, phân bố cụ thể của các phần tử trong tập ảnh hưởng
rất mạnh đến số lượt cần thiết, nên sai số của xấp xỉ này có thể làm chiến
lược kém đi.

Do \(\SetS[a][b]\) nhỏ, về nguyên tắc có thể duyệt toàn bộ tương lai từ tập đó
và dùng cắt tỉa tốt để tính chính xác
\(\operatorname{Step}(\SetS[a][b])\). Bài gốc không hiện thực hướng này vì độ
khó lập trình cao.

Một hướng khác là học một hàm \(\operatorname{Step}\) chính xác hơn theo
\(|\SetS[a][b]|\), chẳng hạn kết hợp các hàm cơ sở như
\(\ln|\SetS[a][b]|\) và \(|\SetS[a][b]|^k\). Ta chạy thử, thống kê với mỗi
kích thước tập ứng viên cần bao nhiêu lượt đoán, dùng thống kê đó xây dựng
hàm \(\operatorname{Step}\) mới, rồi lặp lại. Đây là tinh thần của thuật toán
di truyền.

Tác giả hiện thực một phiên bản đơn giản của quá trình này: mỗi vòng lặp xây
dựng hàm \(\operatorname{Step}\) kế tiếp từ hàm trước. Kết quả thực tế không
quá tốt: sau khoảng \(3\) giờ và hơn \(100\) vòng lặp, chỉ thu được hàm có số
lượt trung bình \(5.250\).

Trong lúc điều chỉnh, tác giả tìm được hàm
\[
  \operatorname{Step}(\SetS[a][b])=\ln(|\SetS[a][b]|+1),
\]
cho số lượt trung bình \(5.243\), tốt hơn mọi chiến lược trước đó. Một lý do
có thể là khi tập còn lại nhỏ, hàm này khớp thực tế hơn.

\subsection{Trò chơi có nhiều vị trí hơn}

Khi số vị trí tăng, số đáp án có thể tăng rất nhanh. Với luật không lặp và
\(10\) vị trí dùng chữ số \(0\) đến \(9\), có \(10! = 3628800\) đáp án.
Khi đó việc tính \(\SetS[a][b]\) trở nên không khả thi, nên các heuristic ở
trên không dùng trực tiếp được.

Một cách xử lý tốt là phân trường hợp. Khi tập ứng viên còn lớn, dùng chiến
lược đơn giản; khi tập đã giảm đủ nhỏ, chuyển sang heuristic hoặc thậm chí
chiến lược tối ưu. Lý do là ở giai đoạn đầu, số lượt đoán khác bản chất khá
ít, lựa chọn khả thi không nhiều, và chất lượng các lượt đoán cũng không khác
nhau quá lớn. Như đã thấy trong phần tối ưu, cách phân đoạn như vậy có thể
tiết kiệm rất nhiều thời gian chạy.

\section{Kết luận}

\subsection{Tổng kết}

Bài viết trình bày phương pháp lọc, thuật toán phổ biến nhất để máy tính giải
trò chơi đoán số, cùng nhiều chiến lược trên nền thuật toán đó: chiến lược đơn
giản, chiến lược heuristic và các thử nghiệm tối ưu hóa. Bảng sau gom các số
liệu chính.

\begin{center}
\small
\begin{tabular}{|l|l|r|r|r|}
\hline
Nhóm & Chiến lược & Nhiều nhất & Trung bình & Thời gian \\
\hline
Đơn giản & Nhỏ nhất theo từ điển & 9 & 5.560 & 10.367 giây \\
Đơn giản & Ngẫu nhiên & 9 & 5.465 & 10.508 giây \\
\hline
Heuristic & Tập tệ nhất & 7 & 5.389 & 79.620 giây \\
Heuristic & Tập trung bình & 7 & 5.268 & 73.762 giây \\
Heuristic & Số bước kỳ vọng & 7 & 5.265 & 95.682 giây \\
Heuristic & Số loại phản hồi & 8 & 5.308 & 78.632 giây \\
\hline
Kết hợp & Thuật toán di truyền & 7 & 5.250 & 97.130 giây \\
Kết hợp & Hàm \(\operatorname{Step}\) tối ưu & 7 & 5.243 & 95.562 giây \\
\hline
\end{tabular}
\end{center}

Thời gian của các chiến lược kết hợp không tính thời gian dùng thuật toán di
truyền hoặc điều chỉnh thủ công để tìm tham số; phần đó trong bài gốc mất
nhiều giờ.

So sánh theo chiều dọc cho thấy heuristic có hiệu quả tổng hợp tốt nhất. Ưu
điểm lớn nhất là hàm đánh giá làm thuật toán có vẻ ``biết suy luận'' hơn: nó
kết hợp khả năng tính toán mạnh của máy tính với một mô hình đánh giá gần với
cách con người cân nhắc thông tin. Ngoài ra, khi tối ưu tính
\(\SetS[a][b]\), cả mục tối ưu lớp tương đương và mục nhiều vị trí đều dùng
cùng một kỹ thuật quan trọng: chia trường hợp. Dù không làm giảm bậc độ phức
tạp trong mọi mô hình, kỹ thuật này cải thiện hiệu quả chạy rất nhiều.

Nhìn chung, giải trò chơi đoán số bằng máy tính có thể xem là một bài toán trí
tuệ nhân tạo nhỏ: câu hỏi cốt lõi là làm sao để máy tính có năng lực suy luận
tốt hơn.

\subsection{Hướng nghiên cứu tiếp theo}

Với luật chuẩn, vẫn còn không gian cải thiện. Chiến lược tối ưu từng phần hoặc
thuật toán di truyền đầy đủ có thể giảm số lượt trung bình ở mức bản chất hơn.
Hiệu quả của các chiến lược này trên những luật khác, đặc biệt các biến thể
Mastermind, cũng rất đáng nghiên cứu.

Ngoài ra, chiến lược đối kháng được nhắc ở phần tiêu chuẩn thống nhất cũng là
một vấn đề lý thuyết trò chơi thú vị: người ra số thay đổi đáp án ngầm nhưng
vẫn giữ mọi phản hồi nhất quán. Bài toán đó thậm chí có thể khó hơn chính bài
toán giải đoán số.

\section{Tài liệu tham khảo}

\begin{enumerate}
  \item Barteld Kooi, University of Groningen. \textit{Yet Another Mastermind
  Strategy}, 2005.
  \item Donald Knuth, Stanford. \textit{The Computer As Master Mind}, 1976.
  \item Một lời giải cho trò chơi đoán số:
  \url{http://www.cppblog.com/lemene/archive/2007/11/26/37273.html}.
  \item Mastermind, Wikipedia:
  \url{http://en.wikipedia.org/wiki/Mastermind_(board_game)}.
  \item Mục từ về trò chơi đoán số trên Baidu Baike.
  \item Lời giải bài \texttt{guess}, Winter Camp 2006.
\end{enumerate}

\appendix
\section{Phụ lục của bài gốc}

Phụ lục bản gốc là phần phản hồi cho môn ``Dẫn luận Khoa học và Công nghệ
Thông tin'', không liên quan trực tiếp đến thuật toán. Tác giả liệt kê ba
giảng viên yêu thích: giáo sư Sun Maosong, viện sĩ Sun Jiaguang và viện sĩ
Zhang Bo. Phần góp ý mong nội dung bài giảng hệ thống hơn để sinh viên hiểu rõ
khung tổng thể của khoa học và công nghệ thông tin; tác giả cũng đề xuất làm
tiểu luận theo nhóm, vì tự nghiên cứu bài toán đoán số có nhiều ý tưởng nhưng
khó hiện thực hết một mình, trong khi năng lực trao đổi và hợp tác nhóm rất
quan trọng trong nghiên cứu.

\end{document}
